% Options for packages loaded elsewhere
\PassOptionsToPackage{unicode}{hyperref}
\PassOptionsToPackage{hyphens}{url}
\PassOptionsToPackage{dvipsnames,svgnames*,x11names*}{xcolor}
%
\documentclass[
  10pt,
]{article}
\usepackage{lmodern}
\usepackage{amssymb,amsmath}
\usepackage{ifxetex,ifluatex}
\ifnum 0\ifxetex 1\fi\ifluatex 1\fi=0 % if pdftex
  \usepackage[T1]{fontenc}
  \usepackage[utf8]{inputenc}
  \usepackage{textcomp} % provide euro and other symbols
\else % if luatex or xetex
  \usepackage{unicode-math}
  \defaultfontfeatures{Scale=MatchLowercase}
  \defaultfontfeatures[\rmfamily]{Ligatures=TeX,Scale=1}
  \setmainfont[]{DejaVu Serif}
  \setmonofont[]{DejaVu Sans Mono}
\fi
% Use upquote if available, for straight quotes in verbatim environments
\IfFileExists{upquote.sty}{\usepackage{upquote}}{}
\IfFileExists{microtype.sty}{% use microtype if available
  \usepackage[]{microtype}
  \UseMicrotypeSet[protrusion]{basicmath} % disable protrusion for tt fonts
}{}
\makeatletter
\@ifundefined{KOMAClassName}{% if non-KOMA class
  \IfFileExists{parskip.sty}{%
    \usepackage{parskip}
  }{% else
    \setlength{\parindent}{0pt}
    \setlength{\parskip}{6pt plus 2pt minus 1pt}}
}{% if KOMA class
  \KOMAoptions{parskip=half}}
\makeatother
\usepackage{xcolor}
\IfFileExists{xurl.sty}{\usepackage{xurl}}{} % add URL line breaks if available
\IfFileExists{bookmark.sty}{\usepackage{bookmark}}{\usepackage{hyperref}}
\hypersetup{
  colorlinks=true,
  linkcolor=Maroon,
  filecolor=Maroon,
  citecolor=Blue,
  urlcolor=Blue,
  pdfcreator={LaTeX via pandoc}}
\urlstyle{same} % disable monospaced font for URLs
\usepackage[margin=0.6in,landscape]{geometry}
\usepackage{longtable,booktabs}
% Correct order of tables after \paragraph or \subparagraph
\usepackage{etoolbox}
\makeatletter
\patchcmd\longtable{\par}{\if@noskipsec\mbox{}\fi\par}{}{}
\makeatother
% Allow footnotes in longtable head/foot
\IfFileExists{footnotehyper.sty}{\usepackage{footnotehyper}}{\usepackage{footnote}}
\makesavenoteenv{longtable}
\setlength{\emergencystretch}{3em} % prevent overfull lines
\providecommand{\tightlist}{%
  \setlength{\itemsep}{0pt}\setlength{\parskip}{0pt}}
\setcounter{secnumdepth}{-\maxdimen} % remove section numbering
\renewcommand{\arraystretch}{1.6}

\author{}
\date{}

\begin{document}

\hypertarget{phase-0.2-keplers-equation-by-hand}{%
\section{Phase 0.2 --- Kepler's Equation by
Hand}\label{phase-0.2-keplers-equation-by-hand}}

\textbf{Type:} Tool drill, no mission narrative. \textbf{Target time:}
\textasciitilde20 minutes per example once the iteration pattern clicks
--- budget longer the first time through. \textbf{Prerequisite:} 0.0
Slide rule basics, 0.1 Trig/log table interpolation.

This is the drill the handoff notes call out specifically: solving
Kepler's equation by hand, iteratively, with a slide rule and trig
tables, is something nobody in the group has actually done --- even
people with graduate-level orbital mechanics training almost always
learned it as code, not as a hand procedure. This is where that changes.

\hypertarget{the-equation}{%
\subsection{The equation}\label{the-equation}}

Kepler's equation relates mean anomaly (time-like, uniform) to eccentric
anomaly (geometric, on the auxiliary circle):

\begin{verbatim}
M = E − e·sin(E)          (radians)
\end{verbatim}

Worked in \textbf{degrees} (the practical choice with a degree-based
trig table), this becomes:

\begin{verbatim}
M = E − k·e·sin(E)          where k = 180/π ≈ 57.29578
\end{verbatim}

There's no closed-form solution for E --- it's solved iteratively.
Newton-Raphson on this equation, in degrees:

\begin{verbatim}
g(E)  = E − k·e·sin(E) − M
g'(E) = 1 − e·cos(E)
E_new = E − g(E)/g'(E)
\end{verbatim}

\textbf{Watch the units carefully here} --- this is the single most
common place to introduce an error.
\texttt{g\textquotesingle{}(E)\ =\ 1\ −\ e·cos(E)} has \textbf{no}
factor of k in it, even though \texttt{g(E)} does. That's not a typo:
differentiating \texttt{sin(E)} with respect to E when E is expressed in
degrees pulls in a factor of π/180 from the chain rule, which exactly
cancels the k in the sine term. If you find yourself writing
\texttt{1\ −\ k·e·cos(E)}, stop --- that's the mistake to check for
first.

A good starting guess: \textbf{E₀ = M}. For higher eccentricity, E₀ = M
+ k·e·sin(M) converges faster (it's literally the first pass of simple
successive substitution) --- worth trying on Example 2 below if you want
to compare iteration counts.

\hypertarget{reference-material}{%
\subsection{Reference material}\label{reference-material}}

Trig table: use the 0.1 excerpt/technique (interpolate sin and cos to at
least 5 decimal places --- errors here feed directly into g(E) and
g'(E)).

Period-authentic note: purpose-built tables existed historically for
solving Kepler's equation directly by lookup rather than iteration
(e.g.~tables of E as a function of M and e, in the tradition of
Bauschinger's orbit-determination tables). Sourcing and linking a
genuine period edition of one of these is flagged as a follow-up for
\texttt{reference/tables.md} --- not yet done, so this drill uses direct
Newton-Raphson iteration instead, which is itself a legitimate
period-appropriate method, just more laborious than a purpose-built
lookup table would have been.

\hypertarget{example-1-low-eccentricity}{%
\subsection{Example 1 --- Low
eccentricity}\label{example-1-low-eccentricity}}

\textbf{Given:} e = 0.100, M = 30.000°

Iterate E\_new = E − g(E)/g'(E) starting from E₀ = M, until
\textbar ΔE\textbar{} is comfortably below your target precision (aim
for \textless{} 0.01°).

\hypertarget{example-2-higher-eccentricity}{%
\subsection{Example 2 --- Higher
eccentricity}\label{example-2-higher-eccentricity}}

\textbf{Given:} e = 0.400, M = 45.000°

Same procedure. Expect more iterations and a bigger first correction ---
this is exactly why higher-eccentricity orbits were the harder case for
real human computers, and why a smarter initial guess mattered more as e
grew.

Once E converges for either example, get the true anomaly ν:

\begin{verbatim}
tan(ν/2) = √[(1+e)/(1−e)] · tan(E/2)
\end{verbatim}

\begin{center}\rule{0.5\linewidth}{0.5pt}\end{center}

\hypertarget{worksheet}{%
\subsection{Worksheet}\label{worksheet}}

Use one row per Newton-Raphson iteration. Add rows as needed.

\textbf{Example 1} (e = 0.100, M = 30.000°)

\begin{longtable}[]{@{}lllllll@{}}
\toprule
n & E\_n (°) & sin(E\_n) & cos(E\_n) & g(E\_n) & g'(E\_n) &
ΔE\tabularnewline
\midrule
\endhead
0 & 30.000 & & & & &\tabularnewline
1 & & & & & &\tabularnewline
2 & & & & & &\tabularnewline
\bottomrule
\end{longtable}

Converged E = \_\_\_\_\_\_\_\_\_\_° → ν = \_\_\_\_\_\_\_\_\_\_°

\textbf{Example 2} (e = 0.400, M = 45.000°)

\begin{longtable}[]{@{}lllllll@{}}
\toprule
n & E\_n (°) & sin(E\_n) & cos(E\_n) & g(E\_n) & g'(E\_n) &
ΔE\tabularnewline
\midrule
\endhead
0 & 45.000 & & & & &\tabularnewline
1 & & & & & &\tabularnewline
2 & & & & & &\tabularnewline
3 & & & & & &\tabularnewline
\bottomrule
\end{longtable}

Converged E = \_\_\_\_\_\_\_\_\_\_° → ν = \_\_\_\_\_\_\_\_\_\_°

\end{document}
